\begin{table}[t]
\centering
\small
\resizebox{\ifdim\width>\linewidth\linewidth\else\width\fi}{!}{%%
\begin{tabular}{llrlrrrlr}
\toprule
prefix & n & prevalence & register levels & base auc & with auc & V & interval & resolved \\
\midrule
0 & 46606 & 0.945 & 3019 & 0.641 & 0.898 & 0.257 & [+0.233, +0.281] & True \\
1 & 46605 & 0.945 & 3019 & 0.733 & 0.900 & 0.166 & [+0.145, +0.203] & True \\
2 & 11080 & 0.722 & 1567 & 0.711 & 0.775 & 0.064 & [+0.047, +0.088] & True \\
3 & 6583 & 0.612 & 1224 & 0.799 & 0.816 & 0.017 & [-0.025, +0.039] & False \\
5 & 3005 & 0.373 & 700 & 0.743 & 0.787 & 0.044 & [+0.021, +0.087] & True \\
8 & 857 & 0.550 & 341 & 0.672 & 0.741 & 0.069 & [-0.001, +0.126] & False \\
\bottomrule
\end{tabular}%
}
\caption{A PREFIX AXIS ON ONE LOG. At prefix $k$ the analyst stands after the case's $k$-th assignment event: the population is the cases still at risk --- more than $k$ events and no group change yet --- the target is whether a change occurs AFTER event $k$, and the baseline is the intake block plus what the prefix has revealed. $k = 0$ is case creation, which is the prediction point every other result in this paper uses, and it is the anchor the rest are read against. Every interval is the pointwise basic construction under the block-weighted bootstrap at \nPrefixDraws\ draws.}
\label{tab:prefix}
\end{table}
